# theorem thm:p148-reconstruction-dichotomy-v5

## statement

Let $X$ be a connected compact Kähler manifold, and let $\theta$ be a smooth
closed real $(1,1)$-form such that $\operatorname{PSH}(X,\theta)\ne\varnothing$.
Write

$$
\mathscr D_{\mathcal I}(X,\theta)
:=\operatorname{PSH}(X,\theta)/{\sim_{\mathcal I}}
$$

for the set of *realizable* $\mathcal I$-data. Thus an element is induced by
an actual $\varphi\in\operatorname{PSH}(X,\theta)$ and is represented by the
tower $(\mathcal I(m\varphi))_{m\ge1}$, equivalently by its divisorial Lelong
profile; it is not an arbitrary formal sequence of ideals. For
$d=[\varphi]_{\mathcal I}$, define

$$
\mathcal R_\theta(d):=P_\theta[\varphi]_{\mathcal I}.
$$

**Theorem (P14.8 v5).**

1. The map $\mathcal R_\theta$ is well defined and returns the unique
   normalized $\mathcal I$-model potential in the class $d$. Hence full
   realizable $\mathcal I$-data reconstructs its canonical normalized
   $\mathcal I$-model representative, with no hypothesis on nonpluripolar
   mass.

2. Raw nonpluripolar mass is not part of $\mathcal I$-data. On
   $X=\mathbf P^1$, equipped with the Fubini--Study form $\omega$ normalized
   by $\int_X\omega=1$, there is one $\mathcal I$-class containing potentials
   $\varphi_s$, $0\le s\le1$, such that

   $$
   [\varphi_s]_{\mathcal I}=[0]_{\mathcal I},\qquad
   \int_X\left\langle\omega+\operatorname{dd}^c\varphi_s\right\rangle=1-s.
   $$

   In particular, the raw zero-mass representative $\varphi_1$ has canonical
   normalized $\mathcal I$-model representative $0$, whose mass is one. So
   neither a raw zero-mass representative nor its mass can be reconstructed
   from its $\mathcal I$-class.

These settle the two defined readings of the former P14.8 wording: complete
realizable $\mathcal I$-data determines the canonical normalized
$\mathcal I$-model representative, while neither a raw representative nor
its raw mass descends to the $\mathcal I$-quotient.

**Scope.** On a smooth projective variety, the separate all-mass order
isomorphism for full realized DXZ--Boucksom--Jonsson objects in a
pseudoeffective real Néron--Severi class remains distinct. The literal
all-semivaluation ideal formula is used only for a $\mathbf Q$-line-bundle
class. This theorem asserts neither a compact-Kähler nor a genuinely
transcendental extension of that result. Realization of arbitrary formal
ideal or valuative data remains P14.1, and the P14.7 terminal-left-limit
corollary is not used here.

The exact source excerpts used by the proof are appended to its frozen V6
packet: Boucksom's multiplier-ideal notes, BEGZ's nonpluripolar-product
paper, and Kellay--Mashreghi, arXiv:1102.5596, formula (3.4), for the
generalized-Cantor capacity criterion.

## proof

Fix $d=[\varphi]_{\mathcal I}$ and set

$$
\mathcal F_d:=\{\psi\in\operatorname{PSH}(X,\theta):
  \psi\le0,\ [\psi]_{\mathcal I}=d\}.
$$

The family is nonempty because $\varphi-\sup_X\varphi$ belongs to it. It is
directed under finite maxima: generic Lelong numbers obey

$$
\nu(\max\{\psi_1,\psi_2\},E)=
\min\{\nu(\psi_1,E),\nu(\psi_2,E)\}
$$

for every prime divisor $E$ over $X$. Boucksom's valuative
multiplier-ideal criterion identifies the full tower with this profile, so a
finite maximum remains in the same class.

Put $U=(\sup_{\psi\in\mathcal F_d}\psi)^*$. Choquet's lemma gives a
sequence $\psi_j\in\mathcal F_d$ with this upper envelope. Replacing it by
$u_j=\max\{\psi_1,\ldots,\psi_j\}$ yields an increasing sequence of members
of $\mathcal F_d$ with $u_j\nearrow U$ almost everywhere. On a model carrying
a prime divisor $E$, choose a very general point of $E$ simultaneously for
the countably many pullbacks of the $u_j$ and for $U$. At that point their
generic Lelong numbers equal their point Lelong numbers. After adding a
smooth local potential, Boucksom's increasing-sequence Lelong theorem gives

$$
\nu(U,E)=\lim_j\nu(u_j,E)=\nu(d,E).
$$

Apply the valuative multiplier-ideal criterion after every positive integral
scaling to obtain $[U]_{\mathcal I}=d$. Thus $U\in\mathcal F_d$ and is its
largest member; the family depends only on $d$, so this is
$\mathcal R_\theta(d)$ and the construction is idempotent. Also
$\sup_X U=0$, since a positive constant could otherwise be added without
altering multiplier ideals. A normalized $\mathcal I$-model potential in the
class is fixed by the same competitor family and hence equals $U$. This
proves the first assertion.

For the second assertion, construct a generalized Cantor set on a fixed
arc of the unit circle, viewed inside $\mathbf P^1$. Normalize the initial
arc length to $\ell_0=1$ (Kellay--Mashreghi use $2\pi$; this fixed rescaling
does not affect convergence of their criterion) and take
$\ell_m=\exp(-4^m)$ for $m\ge1$. The inequalities

$$
2\ell_1<\ell_0,\qquad
2\ell_{m+1}/\ell_m=2\exp(-3\cdot4^m)<1\quad(m\ge1)
$$

make the construction possible. Give a stage-$m$ interval mass $2^{-m}$.
The limiting probability measure $\mu$ is nonatomic. Kellay--Mashreghi,
formula (3.4), states that this generalized Cantor set has zero logarithmic
capacity exactly when $\sum_m2^{-m}\log(1/\ell_m)$ diverges; here its terms
for $m\ge1$ are $2^m$. Thus the resulting compact set $F\subset\mathbf P^1$
has zero logarithmic capacity and is polar (hence pluripolar).

Let $G_p$ be the normalized Green potential with

$$
\omega+\operatorname{dd}^cG_p=\delta_p,\qquad\sup_XG_p=0,
$$

and set $G_\mu(z)=\int_XG_p(z)\,d\mu(p)$. A point outside $F$ has positive
chordal distance from $F$, giving a uniform finite lower bound for the
integrands; hence $G_\mu\not\equiv-\infty$. Fubini gives
$\omega+\operatorname{dd}^cG_\mu=\mu$. For $0\le s\le1$, define
$\varphi_s=sG_\mu$. Then

$$
\omega+\operatorname{dd}^c\varphi_s=(1-s)\omega+s\mu.
$$

On a Riemann surface a Lelong number is the atomic mass of the positive
current. Both summands are nonatomic, so every Lelong number of every
$\varphi_s$ vanishes. As prime divisors over a smooth curve are points, the
valuative criterion at every integral scaling gives
$\mathcal I(m\varphi_s)=\mathcal O_X$ for every $m\ge1$. Therefore
$[\varphi_s]_{\mathcal I}=[0]_{\mathcal I}$.

BEGZ's degree-one nonpluripolar locality applies off the polar set $F$, where
$G_\mu$ is locally bounded, and its product gives no mass to $F$. Hence

$$
\mathbf1_{X\setminus F}\left\langle
  \omega+\operatorname{dd}^c\varphi_s\right\rangle
=(1-s)\mathbf1_{X\setminus F}\omega.
$$

The smooth measure also gives the polar set zero mass; consequently

$$
\left\langle\omega+\operatorname{dd}^c\varphi_s\right\rangle=(1-s)\omega,
\qquad
\int_X\left\langle\omega+\operatorname{dd}^c\varphi_s\right\rangle=1-s.
$$

Applying the first assertion to the common class gives
$\mathcal R_\omega([\varphi_s]_{\mathcal I})=0$. At $s=1$, the input has
mass zero and its canonical normalized $\mathcal I$-model has mass one,
which proves the non-reconstruction assertion.


## Attached immutable source excerpts

The following verbatim excerpts are the only external source interfaces used
in the V6 proof.  They are copied from the immutable local texts named below;
the full-file SHA-256 values make the excerpts reproducible.

### Boucksom — sequential and valuative Lelong facts

Source: `boucksom_l2.txt`
SHA-256: `4a3fa01a21b1cc66fe174460732fe80a29914f26279c04d2824223fc7d9ab0a4`

> Lemma 2.6. The Lelong number `νa(u)` is usc with respect to `a`, and
> with respect to the weak topology of psh functions. In particular, if
> `(uj)` is an increasing sequence of psh functions converging a.e. to a
> psh function `u), then `νa(uj) ↘ νa(u)`.

> Lemma 10.5. We have `νy(φ)=νY(φ)` for a very general point `y∈Y`,
> i.e. outside a countable union of strict closed analytic subsets of `Y`.

> Corollary 10.18. A function `f∈OX` belongs to `J(φ)` iff there exists
> `ε>0` such that `ordE(f) ≥ (1+ε)νE(φ)-AX(E)` for all prime divisors
> `E` over `X`.

These are used respectively for the increasing-envelope limit, simultaneous
very-general-point selection, and the integral-scaled multiplier-ideal
comparison.

### BEGZ — locality and polar-set exclusion

Source: `begz_0812_3674v2.txt`
SHA-256: `890173e309165984432e7326d13b7e096662cc805fbc19045a297f12211254c7`

> Proposition 1.4. Let `u1,...,up` be psh functions on `X`. The
> non-pluripolar product operator is local in the plurifine topology whenever
> well-defined; it depends only on the currents and not on the specific
> potentials.

The same source states in its introduction that a non-pluripolar product
puts no mass on pluripolar subsets.  These two facts yield the degree-one
calculation off the polar Cantor set and then remove the polar set.

### Kellay--Mashreghi — generalized-Cantor capacity criterion

Source: `km_1102_5596.txt`
SHA-256: `fd59393dbbd2c98820bf217b5bd94ec1b786e125c7fa53b00900d56f6bc19cf7`

> The generalized Cantor set has logarithmic capacity zero if and only if
> `sum 2^{-n} log(1/ell_n)=∞` (formula (3.4)).

For `ell_m=exp(-4^m)`, the terms are `2^m`, so the series diverges.  The
proof separately checks the child-fitting inequalities before applying this
criterion.

